\infographicchapter{Mammography}{mamography.png}
\label{ch:infographic-mamography}

\section{Overview}

Mammography is a dedicated, low-dose X-ray exam designed specifically to look inside breast tissue. Because fatty tissue and milk-gland tissue in the breast look fairly similar to regular X-rays, mammography machines use specially tuned low-energy X-ray beams to bring out subtle differences. Gently but firmly pressing the breast between two smooth plates spreads the tissue out evenly, making the picture sharper, keeping it still, and lowering the amount of radiation needed.

Doctors use mammography in two main ways. A screening mammogram checks people who have no symptoms, aiming to find breast cancer very early---often years before a lump can be felt by hand. A diagnostic mammogram takes extra, closer views when someone has a symptom like a lump, localized pain, skin thickening, or nipple discharge, or when a screening exam shows something that needs a second look.

Mammography is the proven gold standard for breast cancer screening because it saves lives, takes only a few minutes, and is widely available. Its main limits are that dense breast tissue (which looks white on the scan, just like many cancers) can hide small tumors, and overlapping tissue on flat 2D pictures can sometimes trigger callbacks for extra views. Modern 3D mammography (digital breast tomosynthesis) helps reduce this overlap. Digital mammography uses a very small X-ray dose---about 0.4 millisieverts (mSv) for a standard two-view exam, roughly equal to a few weeks of everyday natural background radiation---so the benefits of early detection strongly outweigh the tiny theoretical risk.

\section{Underlying Physical Principles}

\begin{itemize}
    \item \textbf{Specially Tuned Low-Energy X-Rays:} Unlike general X-ray tubes that use tungsten, mammography machines use special molybdenum or rhodium metal targets and filters to produce a narrow beam of gentle, low-energy X-rays. This tuned beam highlights subtle differences inside soft breast tissue while filtering out rays that would add dose without improving the image.
    \item \textbf{High Soft-Tissue and Calcium Contrast:} Low-energy X-rays are absorbed differently by fat, glandular tissue, and tiny specks of calcium (microcalcifications). This makes dense soft-tissue lumps stand out and causes tiny calcium deposits to appear crisp and bright white.
    \item \textbf{Gentle Firm Compression and Sharp Detectors:} Flattening the breast evenly reduces tissue thickness, prevents motion blur, cuts down scattered X-rays, and lowers radiation dose. X-rays passing through the breast hit a high-resolution digital flat-panel sensor that captures tiny details.
    \item \textbf{2D Views, 3D Tomosynthesis, and Contrast Views:} Standard digital mammograms create detailed flat 2D pictures. In 3D mammography (digital breast tomosynthesis), the X-ray arm moves in a short arc to take several low-dose angled pictures that a computer builds into thin 1-millimeter slices. Contrast-enhanced mammography adds an IV dye to highlight areas of extra blood flow.
\end{itemize}

\section{Imaging Mechanism}

\begin{enumerate}
    \item \textbf{Positioning and Compression:} A technologist gently places one breast at a time onto the flat detector platform and lowers a smooth plastic paddle to flatten and hold the tissue steady for two standard angles: top-to-bottom (CC) and side-angled (MLO).
    \item \textbf{Automatic Low-Dose X-Ray Exposure:} Sensors measure the thickness and density of the breast and automatically choose the lowest X-ray setting needed for a clear picture. For 3D tomosynthesis, the X-ray tube glides in a smooth arc over a few seconds to capture multiple angled views.
    \item \textbf{High-Resolution Digital Capture:} X-rays passing through the breast travel through a grid that blocks stray scatter and strike a fine-detail digital detector, turning the pattern into crisp electronic image data.
    \item \textbf{Computer Processing and 3D Slicing:} Specialized software balances brightness across the whole breast and sharpens fine structures. For 3D exams, the computer reconstructs thin 1-millimeter slices so the radiologist can scroll through the breast layer by layer.
    \item \textbf{Radiologist Review:} The finished images are displayed on high-brightness medical monitors side by side with any previous mammograms so the radiologist can spot even the smallest changes over time.
\end{enumerate}

\section{Key Technical Parameters}

\noindent\begin{tabularx}{\textwidth}{@{}l X@{}}
\toprule
\textbf{Parameter} & \textbf{Typical Value / Description} \\
\midrule
Physical Principle & Low-energy X-rays pass through gently compressed breast tissue to highlight differences in fat, glandular tissue, and tiny calcium specks \\
Energy / Signal Source & Dedicated mammography X-ray tube with specialized filters and a high-resolution digital detector \\
Typical Frequency / Energy & Low-energy X-ray beam (25--35 kVp) \\
Spatial Resolution & Extremely fine detail (50--100 micrometer pixels, sharp enough to see tiny microcalcifications) \\
Typical Scan Duration & 10--15 minutes for routine screening; 20--30 minutes for diagnostic or extra magnified views \\
Contrast Agent & None for standard 2D or 3D mammography; IV iodinated contrast dye for contrast-enhanced mammography (CEM) \\
Ionizing Radiation & Yes, very low dose (about 0.4 mSv for a standard two-view exam, well within strict safety limits) \\
\bottomrule
\end{tabularx}

\section{Common Symptoms That Prompt This Scan}
\begin{itemize}
    \item \textbf{Routine Screening With No Symptoms:} Most mammograms are performed as routine preventive screening in women with no breast complaints to catch cancer before it can be felt.
    \item \textbf{A New Breast Lump or Thickening:} A firm lump, thickened patch, or new fullness in the breast or underarm noticed by you or your doctor.
    \item \textbf{Skin or Nipple Changes:} Dimpling, puckering, redness, scaly skin around the nipple, a newly inverted nipple, or spontaneous bloody or clear nipple discharge.
    \item \textbf{Localized, Persistent Breast Pain:} Pain focused in one specific spot of the breast that does not come and go with normal menstrual cycles.
\end{itemize}

\section{Clinical Usages}

\subsection{Primary Usages}

\begin{itemize}
    \item \textbf{Routine Breast Cancer Screening:} Regular checkups for people without symptoms (using 2D digital mammography or 3D tomosynthesis) to catch early invasive cancer and pre-invasive milk-duct cancer (DCIS) before a lump can be felt---recommended every 2 years from ages 40 to 74 for women at average risk, or earlier for those at higher risk.
    \item \textbf{Checking New Breast Symptoms:} First-choice imaging for women aged 30 and older (and selected men) who notice a breast lump, skin thickening, nipple pulling inward, or unusual nipple discharge.
    \item \textbf{Follow-Up of Screening Findings:} Extra magnified or focused compression pictures to take a closer look at any uncertain area, uneven density, or calcium cluster spotted on a screening mammogram.
    \item \textbf{Spotting Microcalcifications:} The best test for finding and evaluating tiny, clustered calcium specks that can be the earliest clue to ductal carcinoma in situ (DCIS) or early breast cancer.
    \item \textbf{Guiding Biopsies and Surgery:} Helps guide needle biopsies (stereotactic or 3D-guided) of tiny calcium clusters or tissue distortion, and helps surgeons mark non-palpable spots accurately before breast-conserving surgery.
\end{itemize}

\subsection{Secondary Usages}

\begin{itemize}
    \item \textbf{Monitoring After Breast Cancer Treatment:} Yearly follow-up mammograms after lumpectomy or single-side mastectomy to watch for returning cancer or a new cancer in the other breast.
    \item \textbf{Checking Surgical Tissue Samples:} X-raying removed breast tissue right in the operating room to make sure the entire targeted spot, calcium cluster, and marker clip were completely removed.
    \item \textbf{Contrast-Enhanced Mammography (CEM):} Uses an IV contrast dye to highlight increased blood flow in dense breasts, measure tumor extent before surgery, or check response to chemotherapy when breast MRI is not an option.
    \item \textbf{High-Risk Screening Partner:} Paired with yearly breast MRI in people with a high lifetime risk of breast cancer, such as BRCA1/2 gene mutation carriers or people who had chest radiation earlier in life.
    \item \textbf{Galactography (Ductography) for Nipple Discharge:} A specialized contrast mammogram in which a tiny drop of dye is gently placed into a milk duct to investigate spontaneous, unexplained clear or bloody nipple discharge and spot small intraductal growths or early ductal cancer.
\end{itemize}

\section{Advantages}

\begin{itemize}
    \item \textbf{Proven to Save Lives:} Mammography is the only breast imaging test proven in large clinical trials to lower the risk of dying from breast cancer by catching disease early.
    \item \textbf{Unmatched for Tiny Calcium Specks:} Finds and characterizes fine microcalcifications better than any other test, often spotting early ductal cancer (DCIS) before a mass forms.
    \item \textbf{Very Low Radiation and Fast Exam:} Uses a tiny X-ray dose (about 0.4 mSv for a standard exam) and takes only 10 to 15 minutes from start to finish.
    \item \textbf{Widely Available and Standardized:} Offered at clinics and hospitals almost everywhere, using a clear, standardized rating system (BI-RADS) that makes year-to-year comparisons reliable.
    \item \textbf{3D Slicing and Biopsy Guidance:} 3D tomosynthesis reduces overlapping tissue shadows and improves cancer detection, while stereotactic guidance allows accurate, minimally invasive needle biopsies.
\end{itemize}

\section{Disadvantages}

\begin{itemize}
    \item \textbf{Harder to Read in Dense Breasts:} Dense glandular tissue appears bright white on a mammogram, just like many cancers, which can hide tumors and make extra tests like ultrasound or MRI helpful.
    \item \textbf{Callbacks and False Alarms:} Overlapping normal tissue on flat 2D views can look like a lump (3D tomosynthesis helps reduce this), sometimes leading to extra follow-up visits, anxiety, or biopsies that turn out to be benign.
    \item \textbf{Brief Compression Discomfort:} Firm pressure between the paddles is necessary for sharp, low-dose pictures, but it can feel tight or uncomfortable for a few seconds per view.
    \item \textbf{Cannot Always Tell Cysts from Solid Lumps:} A mammogram alone cannot always distinguish a harmless fluid-filled cyst from a solid mass without a quick breast ultrasound, and implants or old scar tissue can partly hide nearby tissue.
    \item \textbf{Uses Low-Dose X-Rays:} Because mammography uses X-rays, doctors weigh benefits carefully during pregnancy and generally start with radiation-free ultrasound for women under age 30 who have a new lump.
\end{itemize}

\section{Contraindications}

\subsection{Absolute}

\begin{itemize}
    \item \textbf{None for Standard Mammography:} There are no strict absolute reasons to forbid a standard mammogram when it is medically needed (such as checking a suspicious lump), because the X-ray dose is very low.
\end{itemize}

\subsection{Relative}

\begin{itemize}
    \item \textbf{Pregnancy, Breastfeeding, and Young Age:} Routine screening mammograms are usually paused during pregnancy. For pregnant or breastfeeding patients and women under age 30 with a new symptom, targeted breast ultrasound is preferred first, saving mammography for suspicious findings.
    \item \textbf{Local Skin Issues and Contrast Precautions:} Severe breast infection (mastitis), open skin wounds, or very recent unhealed breast surgery may make compression too painful or unsafe. If contrast-enhanced mammography (CEM) is planned, severe iodinated contrast allergy or severe kidney disease must be considered.
\end{itemize}

\section{Patient Safety Guidelines}

\begin{itemize}
    \item \textbf{Accredited Equipment and Lowest Safe Dose:} Always have your mammogram at a certified facility (accredited under FDA MQSA standards in the United States), where physicists regularly test machines to keep radiation as low as possible while maintaining crisp image quality.
    \item \textbf{Special Care for Breast Implants:} Women with cosmetic breast implants should continue routine screening, while most people who had implant reconstruction after a full mastectomy usually do not need a mammogram of the reconstructed breast. Tell the technologist about implants before positioning so special implant-displacement views can use gentle, carefully controlled compression to protect the implant while clearly showing your natural breast tissue.
    \item \textbf{Pregnancy and Age-Appropriate Choices:} Staff always ask about possible pregnancy or breastfeeding. Ultrasound is usually the first step for pregnant patients or symptomatic people under 30, with mammography added only when truly needed.
    \item \textbf{Comfortable, Safe Compression:} The technologist checks for fragile skin, wounds, or recent incisions and communicates with you while lowering the paddle so pressure is firm enough for clear images without causing unnecessary pain.
    \item \textbf{Clear Reporting and Density Notice:} Facilities must send timely, standardized BI-RADS reports to your doctor and send you a clear written summary of your results and breast density so any needed follow-up happens promptly.
\end{itemize}

\section{Prerequisite Tests}

\begin{itemize}
    \item \textbf{No Routine Blood Tests Needed:} Standard 2D and 3D mammograms require no blood work. Only contrast-enhanced mammography (CEM) may require a kidney function blood test and contrast allergy check beforehand.
    \item \textbf{Previous Mammograms for Comparison:} Bringing or transferring your past mammograms is one of the most helpful steps you can take. Comparing new and old pictures side by side helps the radiologist spot subtle new changes and avoid unnecessary callbacks.
    \item \textbf{Health History and Symptom Check:} The clinic records your personal and family history of breast cancer, prior breast surgeries or implants, and any current symptoms (like lumps or discharge) to choose the right type of exam.
    \item \textbf{Pregnancy and Breastfeeding Check:} Staff confirm whether you might be pregnant or are currently breastfeeding so they can time routine screening safely or start with ultrasound if appropriate.
\end{itemize}

\section{Preparing for the Scan}

\begin{itemize}
    \item \textbf{Skip Deodorant, Powders, and Lotions:} On the day of your exam, do not put deodorant, antiperspirant, body powder, lotion, cream, or perfume on your breasts or underarms. Tiny metal or talc particles in these products can look like suspicious calcium specks on the X-ray.
    \item \textbf{Wear a Two-Piece Outfit:} Choose a separate top and pants or skirt so you only need to undress from the waist up, and take off necklaces or chest jewelry before the pictures.
    \item \textbf{Time Your Visit for Comfort:} If your breasts tend to feel tender before your period, try scheduling your screening mammogram 1 to 2 weeks after your period starts, when breast tissue is least sensitive. For contrast-enhanced mammography (CEM), expect a quick arm IV beforehand.
    \item \textbf{Bring Prior Records and Share Details:} Bring any past mammograms done at another clinic, and tell the technologist before starting if you have breast implants, past surgeries, a new lump, nipple discharge, or could be pregnant or breastfeeding.
\end{itemize}

\section{During Scan}

\begin{itemize}
    \item \textbf{Positioning and Gentle Compression:} You will stand (or sit if needed) in front of the machine while the technologist positions one breast at a time onto the support plate and lowers a smooth plastic paddle to press the breast firmly for two standard pictures on each side.
    \item \textbf{Holding Still and Brief Breath-Holds:} You will be asked to stay completely still and hold your breath for just 2 to 5 seconds as each picture is taken so the image stays sharp and does not need to be repeated.
    \item \textbf{What It Feels and Sounds Like:} The compression feels snug and briefly uncomfortable for a few seconds per view, then releases automatically the moment the picture finishes. You will hear a soft beep or gentle whirring sound, especially during 3D tomosynthesis sweeps.
    \item \textbf{When to Speak Up:} Tell the technologist right away if you feel sharp pain, skin pinching, or lightheadedness so they can adjust your position gently. If you are having a contrast-enhanced mammogram, report any itching, rash, or trouble breathing after the IV dye immediately.
    \item \textbf{How Long It Takes:} A routine screening mammogram usually takes 10 to 15 minutes from start to finish (with actual compression lasting only seconds per picture). A diagnostic exam with extra views typically takes 20 to 30 minutes.
\end{itemize}

\section{Normal Appearance}

\begin{itemize}
    \item \textbf{Balanced Tissue Pattern:} Both breasts look broadly symmetrical from side to side. Radiologists also note your breast density on a standard scale from A (mostly fatty) to D (extremely dense), which is a normal personal trait.
    \item \textbf{Smooth Skin and Clear Outline:} The skin edge looks thin and smooth (usually under 2 millimeters), the nipples point outward normally, and the dark fatty layer behind the breast and chest muscle edge appear clean and clear.
    \item \textbf{Healthy Lymph Nodes and Structures:} Normal underarm lymph nodes look small and oval with a dark fatty center, supporting ligaments look thin, and there are no jagged lumps, suspicious calcium clusters, or pulled-in tissue patterns.
\end{itemize}

\section{Observing for Abnormal Patterns}

\begin{itemize}
    \item \textbf{Irregular or Spiky Lumps:} A dense, uneven lump with radiating spiky edges is the classic sign of invasive breast cancer. In contrast, a smooth, round or oval spot is much more likely to be a harmless fluid cyst or benign fibroadenoma, though ultrasound is used to confirm.
    \item \textbf{Suspicious Microcalcifications:} Clusters of very tiny, irregular, or branching calcium specks lined up along a milk duct can signal early ductal cancer (DCIS) or invasive cancer. By contrast, large coarse ``popcorn'' spots, smooth eggshell rings, or parallel blood-vessel calcium lines are usually harmless.
    \item \textbf{Architectural Distortion:} Thin lines pulling inward toward a single point---even without a clear lump---can be an important clue to invasive breast cancer, though benign radial scars or old surgical scars can look similar.
    \item \textbf{New Asymmetries and Skin Changes:} A new or growing area of uneven tissue compared with past mammograms always gets a closer look. Other warning signs include thickened skin, a newly pulled-in nipple, or enlarged, dense underarm lymph nodes.
    \item \textbf{Harmless Look-Alikes:} Normal overlapping tissue on 2D pictures can mimic a lump until flattened out on extra views or 3D slices. Deodorant specks in the underarm can mimic suspicious calcium, and skin moles or surgical clips can sometimes project over breast tissue.
\end{itemize}

\section{Dealing with Uncertainties}

\begin{itemize}
    \item \textbf{Dense Tissue and Overlapping Shadows:} Dense white breast tissue can hide small cancers or stack together to look like a fake lump. 3D tomosynthesis and focused spot-compression views help separate true findings from normal overlap.
    \item \textbf{Findings That Overlap Between Benign and Cancerous:} Smooth lumps (cysts versus solid tumors), faint calcium clusters, and subtle tissue pulling (old scars versus cancer) cannot always be sorted out by mammography alone.
    \item \textbf{Motion or Skin-Product Artifacts:} Slight movement, incomplete positioning, or leftover deodorant and powder on the skin can blur details or mimic calcium specks, sometimes requiring a quick repeat picture after wiping the skin clean.
    \item \textbf{Clear Step-by-Step Resolution:} If a screening exam is marked incomplete (BI-RADS 0), do not panic---most callbacks turn out to be normal or benign. Extra magnified views, targeted ultrasound, breast MRI, or contrast mammography clarify the picture, and a quick needle biopsy provides a definite answer if anything remains suspicious (BI-RADS 4 or 5).
\end{itemize}

\section{Augmenting Methods}

\begin{itemize}
    \item \textbf{3D Mammography (Digital Breast Tomosynthesis):} Reconstructs thin 1-millimeter slices through the breast to reduce overlapping tissue shadows, find more invasive cancers, and lower false-alarm callbacks.
    \item \textbf{Targeted Breast Ultrasound:} A painless, radiation-free partner test that easily tells harmless fluid cysts from solid lumps, looks through dense breast tissue, checks palpable lumps, and guides needle biopsies in real time.
    \item \textbf{Breast MRI and Contrast-Enhanced Mammography:} Highly sensitive dye-based exams used alongside mammography for people at high lifetime risk (such as BRCA1/2 carriers), very dense breasts, surgery planning, or checking chemotherapy response.
    \item \textbf{Needle Biopsy and AI Support:} Image-guided core needle biopsy gives a definite lab answer and checks tumor hormone markers, while FDA-cleared artificial intelligence tools help radiologists double-check for subtle lumps and calcium specks.
\end{itemize}

\section{Comparison with Alternative Modalities}

\begin{itemize}
    \item \textbf{Compared with Breast Ultrasound:} Ultrasound uses zero radiation and is great for telling cysts from solid lumps and checking dense breasts or women under 30. However, mammography is far better at spotting tiny calcium specks (early DCIS) and is proven to lower breast cancer deaths.
    \item \textbf{Compared with Breast MRI:} Breast MRI is the most sensitive test overall---especially in very dense breasts and high-risk patients---but mammography is much faster, less expensive, widely available, and better at evaluating tiny calcium deposits.
    \item \textbf{Compared with 3D Tomosynthesis (DBT):} 3D tomosynthesis builds thin slices from angled low-dose views to reduce tissue overlap and callbacks compared with standard 2D mammography, using a very similar low radiation dose.
\end{itemize}

\section{Follow-up Steps}

\begin{itemize}
    \item \textbf{Immediate Return to Normal Day and Written Results:} Standard mammography requires zero recovery time. You will usually receive a clear written letter or electronic portal report within about 2 weeks explaining your results and your breast density category, along with whether extra screening might help for dense breasts; contact your provider if your results have not arrived by then.
    \item \textbf{Extra Views for Incomplete Results (BI-RADS 0):} If your screening exam needs a closer look, your clinic will schedule focused compression views, magnified pictures, 3D slices, or a targeted ultrasound to give you a clear answer.
    \item \textbf{Biopsy and Specialist Care (BI-RADS 4--5):} If a spot looks suspicious, a minimally invasive needle biopsy (often with a tiny harmless metal marker clip left to note the spot) confirms what the tissue is. If cancer or a high-risk finding is diagnosed, a coordinated breast care team guides next steps.
    \item \textbf{Ongoing Screening Schedule (BI-RADS 1--3):} Normal or clearly benign results (BI-RADS 1 or 2) simply return to routine screening (usually every 1 to 2 years). Probably benign spots (BI-RADS 3) are typically rechecked in 6 months to confirm they stay stable, and high-risk individuals often pair mammography with yearly breast MRI.
\end{itemize}

\section{Key References}

\begin{itemize}
    \item American College of Radiology (ACR). \textit{ACR BI-RADS Atlas: Breast Imaging Reporting and Data System}, 5th Edition. Reston, VA: ACR, 2013.
    \item American College of Radiology (ACR). \textit{ACR Appropriateness Criteria: Breast Cancer Screening and Palpable Breast Masses}. Reston, VA: ACR, 2023--2024.
    \item U.S. Food and Drug Administration (FDA). \textit{Mammography Quality Standards Act (MQSA) Regulations (21 CFR Part 900)}. Updated 2024.
    \item Monticciolo DL, Malak SF, Friedewald SM, et al. Breast Cancer Screening Recommendations Inclusive of All Women at Average Risk: Update from the ACR and Society of Breast Imaging. \textit{Journal of the American College of Radiology}, 2021;18(9):1280--1288.
    \item National Cancer Institute (NCI). \textit{Mammograms}. National Institutes of Health. Available at: \url{https://www.cancer.gov/types/breast/screening/mammograms}.
    \item Radiopaedia.org. \textit{Mammography}. Radiopaedia. Available at: \url{https://radiopaedia.org/articles/mammography}.
\end{itemize}

\section{Further Reading}
\begin{itemize}
  \item Wikipedia contributors. \textit{Mammography}. Wikipedia, The Free Encyclopedia. Available at: \url{https://en.wikipedia.org/wiki/Mammography}.
\end{itemize}
